%% toc: P15 166 166.4
\chapter{Relative Information and Kullback--Leibler Divergence}
\label{sec:p15-w02b-p15-045-relative-information}

Two hypotheses about the population that produced an observation are
represented by two probability measures on one sample space. Kullback
and Leibler measure how well an observation separates them by the
logarithm of the ratio of the two densities, and average that logarithm
over the first hypothesis. Their 1951 paper extends Shannon's finite
definition to an arbitrary measurable space and shows that no
measurable reduction of the data can increase this average. They also
show that a reduction preserves it exactly when the reduction is a
sufficient statistic \cite[p.~79]{KullbackLeibler1951Information}.

\emph{Purpose.} To compare two fully specified probability models by the
expected discriminating information of one observation, and to decide
whether a statistic loses any of it.
\emph{Inputs.} A measurable space; two probability measures that have the
same null sets; a probability measure with those null sets; the densities;
and, for data reduction, a measurable map into a second measurable space.
\emph{Conclusions.} The information quantities \(I_{1:2}(E)\), \(I(1:2)\)
and \(J(1,2)\) as well-defined elements of \((-\infty,\infty]\) or
\([0,\infty]\); nonnegativity with its equality case; a lower bound on
the information conditional on an event; additivity over independent
components; the data-reduction inequality with its exact equality
condition; the local relation to Fisher's information matrix and the
divergence of two normal populations with a common covariance matrix;
values for pairs whose null sets differ; and a precise list of what these
quantities are not.
Natural logarithms are used unless a base is stated; another base
multiplies every quantity by a positive constant.

\section{The compared measures and the information quantities}
\label{sec:p15-w02b-p15-045-setting}

\begin{hbdefinition}{def:p15-w02b-p15-045-pair}{Compared measures, dominating measure and densities}
Let \((X,\mathbf S)\) be a measurable space and let \(\mu_1,\mu_2\) be
probability measures on \(\mathbf S\). Write \(\mu_1\equiv\mu_2\) when
\(\mu_1\) and \(\mu_2\) have the same null sets, that is, each is
absolutely continuous with respect to the other. In this chapter
\(\equiv\) denotes this relation between measures and nothing else; for
a measure \(\lambda\) and the pair, \(\lambda\equiv\{\mu_1,\mu_2\}\)
means that \(\lambda(E)=0\) exactly when \(\mu_1(E)=0\) and
\(\mu_2(E)=0\).

The \emph{standing hypotheses} are:
\begin{enumerate}
\item \(\mu_1\equiv\mu_2\);
\item \(\lambda\) is a probability measure on \(\mathbf S\) with
\(\lambda\equiv\{\mu_1,\mu_2\}\); for instance \(\lambda=\mu_1\),
\(\lambda=\mu_2\) or \(\lambda=\tfrac12(\mu_1+\mu_2)\);
\item \(f_i\) is a version of the Radon--Nikodym density
\(d\mu_i/d\lambda\) (\cref{thm:p15-radon-nikodym}), so that
\[
  \mu_i(E)=\int_E f_i(x)\,d\lambda(x)\qquad(E\in\mathbf S,\ i=1,2),
\]
and \(0<f_i(x)<\infty\) \([\lambda]\), where \([\lambda]\) means
``except on a \(\lambda\)-null set''.
\end{enumerate}
\(H_i\) is the hypothesis that \(x\) was drawn from the population whose
probability measure is \(\mu_i\).
\end{hbdefinition}

These are the hypotheses of \cite[pp.~79--80, eq.~(2.1)]{KullbackLeibler1951Information}.
Under (1) and (2), densities satisfying (3) always exist: \(\mu_i\ll\lambda\)
gives a density, and \(\mu_i(\{f_i=0\})=0\) together with
\(\lambda\equiv\{\mu_1,\mu_2\}\) gives \(\lambda(\{f_i=0\})=0\).
Kullback and Leibler impose (1) to exclude the case in which some event
has positive probability under one hypothesis and probability zero
under the other; one observation in such an event discriminates
perfectly \cite[p.~79, n.~1]{KullbackLeibler1951Information}. That case
is treated separately in \cref{sec:p15-w02b-p15-045-cases}.

\begin{hbdefinition}{def:p15-w02b-p15-045-information}{Information in an observation, mean information and divergence}
Under the standing hypotheses of \cref{def:p15-w02b-p15-045-pair}:
\begin{enumerate}
\item the \emph{information in \(x\) for discrimination between \(H_1\)
and \(H_2\)} is
\[
  \log\frac{f_1(x)}{f_2(x)};
\]
\item for \(E\in\mathbf S\), the \emph{mean information for
discrimination per observation from \(E\) under \(\mu_1\)} is
\[
  I_{1:2}(E)=\frac1{\mu_1(E)}\int_E d\mu_1(x)\log\frac{f_1(x)}{f_2(x)}
            =\frac1{\mu_1(E)}\int_E f_1(x)\log\frac{f_1(x)}{f_2(x)}\,d\lambda(x)
\]
if \(\mu_1(E)>0\), and \(I_{1:2}(E)=0\) if \(\mu_1(E)=0\);
\item \(I(1:2)=I_{1:2}(X)\);
\item \(I_{2:1}(E)\) and \(I(2:1)\) are defined by exchanging the indices
\(1\) and \(2\);
\item \(J_{12}(E)=I_{1:2}(E)+I_{2:1}(E)\), and the \emph{divergence} is
\[
  J(1,2)=J_{12}(X)=\int\bigl(f_1(x)-f_2(x)\bigr)\log\frac{f_1(x)}{f_2(x)}\,d\lambda(x).
\]
\end{enumerate}
\end{hbdefinition}

These are equations (2.2)--(2.6) of
\cite[pp.~80--81]{KullbackLeibler1951Information}. Two conventions are
part of the definition. First, \(I_{1:2}(E)\) may equal \(+\infty\); the
next lemma shows that it is never \(-\infty\) and never undefined.
Second, the values do not depend on which dominating measure is chosen,
so the restriction of \(\lambda\) to probability measures costs nothing.

\begin{hblemma}{thm:p15-w02b-p15-045-well-defined}{proved=annex:p15-w02b-p15-045-well-defined}{The information integrals exist and do not depend on the dominating measure}
Assume the standing hypotheses of \cref{def:p15-w02b-p15-045-pair}.
\begin{enumerate}
\item For every \(E\in\mathbf S\), the negative part of
\(f_1\log(f_1/f_2)\) has \(\lambda\)-integral over \(E\) at most
\(\mu_2(E)/e\). Hence \(\int_Ef_1\log(f_1/f_2)\,d\lambda\) exists in
\((-\infty,\infty]\), and \(I_{1:2}(E)\in(-\infty,\infty]\). The same
holds with the indices exchanged.
\item \(f_1/f_2\) is a version of \(d\mu_1/d\mu_2\). Consequently, if
\(\lambda'\) is any \(\sigma\)-finite measure on \(\mathbf S\) with
\(\lambda'\equiv\{\mu_1,\mu_2\}\) and \(f_i'=d\mu_i/d\lambda'\), then
\(f_1'/f_2'=f_1/f_2\) \([\lambda]\), and every quantity of
\cref{def:p15-w02b-p15-045-information} computed from
\((\lambda',f_1',f_2')\) equals the one computed from
\((\lambda,f_1,f_2)\).
\item If \(\mu_1(E)>0\), then \(\mu_2(E)>0\) and
\[
  J_{12}(E)=\int_E\Bigl(\frac{f_1(x)}{\mu_1(E)}-\frac{f_2(x)}{\mu_2(E)}\Bigr)
             \log\frac{f_1(x)}{f_2(x)}\,d\lambda(x)\in(-\infty,\infty],
\]
the integral existing in the extended sense. The integrand of
\(J(1,2)\) is nonnegative, \(J(1,2)\in[0,\infty]\), and
\(J(1,2)=I(1:2)+I(2:1)\).
\item If \(E_1,\dots,E_n\in\mathbf S\) are disjoint with union \(X\),
then
\[
  I(1:2)=\sum_{j=1}^n\mu_1(E_j)\,I_{1:2}(E_j),
\]
a sum of terms in \((-\infty,\infty]\), with a term equal to \(0\)
when \(\mu_1(E_j)=0\).
\end{enumerate}
\end{hblemma}

Part (1) is the source's remark that the integral always exists,
possibly as \(+\infty\), because the measures are finite
\cite[p.~80]{KullbackLeibler1951Information}. Part (2) is what makes
Lebesgue measure usable as a dominating measure in the examples below,
although it is not a probability measure.

\begin{hbconvention}{def:p15-w02b-p15-045-conventions}{Support and finiteness conventions}
Every theorem of this chapter states which of the following it assumes.
\begin{enumerate}
\item \emph{Mutual absolute continuity.} Unless a statement says
otherwise, the standing hypotheses of \cref{def:p15-w02b-p15-045-pair}
hold, so \(f_1/f_2\) is finite and positive \([\lambda]\).
\item \emph{Finiteness.} \(I(1:2)\), \(I(2:1)\) and \(J(1,2)\) are
elements of \([0,\infty]\) (\cref{thm:p15-w02b-p15-045-nonnegative});
\(+\infty\) is a legitimate value and is not excluded by the standing
hypotheses. Where a conclusion needs \(I(1:2)<\infty\), the
hypothesis is written in the statement.
\item \emph{Support mismatch.} When \(\mu_1\) and \(\mu_2\) do not have
the same null sets, the standing hypotheses fail and
\cref{def:p15-w02b-p15-045-information} does not apply. The values used
are then those of \cref{def:p15-w02b-p15-045-extension}, and every
statement that uses them says so. That extension covers both one-sided
cases and the case in which each measure charges a set that the other
does not. It is not part of the source; its two clauses are justified by
\cref{thm:p15-w02b-p15-045-support}(1) and~(2).
\end{enumerate}
\end{hbconvention}

\begin{hbdefinition}{def:p15-w02b-p15-045-extension}{Information for a pair whose null sets may differ}
Let \(\mu_1,\mu_2\) be any probability measures on \((X,\mathbf S)\).
\begin{enumerate}
\item If \(\mu_1\ll\mu_2\), let \(g\) be a version of
\(d\mu_1/d\mu_2\) (\cref{thm:p15-radon-nikodym}) with
\(0\leq g<\infty\), and put
\[
  I(1:2)=\int\log g\,d\mu_1,
\]
the extended integral of \cref{thm:p15-w02b-p15-045-extension}(1).
\item If \(\mu_1\) is not absolutely continuous with respect to
\(\mu_2\), that is, \(\mu_1(E)>0=\mu_2(E)\) for some \(E\in\mathbf S\),
put \(I(1:2)=+\infty\).
\item \(I(2:1)\) is defined by exchanging the indices, and
\(J(1,2)=I(1:2)+I(2:1)\in[0,\infty]\).
\end{enumerate}
\end{hbdefinition}

\begin{hblemma}{thm:p15-w02b-p15-045-extension}{proved=annex:p15-w02b-p15-045-extension}{The extended values are well defined and agree with the source's}
Let \(\mu_1,\mu_2\) be probability measures on \((X,\mathbf S)\), and
put \(\varphi(t)=t\log t\) for \(t>0\) and \(\varphi(0)=0\).
\begin{enumerate}
\item If \(\mu_1\ll\mu_2\) and \(g\) is as in
\cref{def:p15-w02b-p15-045-extension}(1), then
\(\mu_1(\{g=0\})=0\), the negative part of \(\log g\) has
\(\mu_1\)-integral at most \(1/e\), and
\[
  I(1:2)=\int\varphi(g)\,d\mu_2\in[0,\infty].
\]
The value does not depend on the version of \(g\), and \(I(1:2)=0\) if
and only if \(\mu_1=\mu_2\).
\item If \(\mu_1\equiv\mu_2\) and \((\lambda,f_1,f_2)\) satisfy the
standing hypotheses of \cref{def:p15-w02b-p15-045-pair}, then
\(I(1:2)\), \(I(2:1)\) and \(J(1,2)\) of
\cref{def:p15-w02b-p15-045-extension} equal those of
\cref{def:p15-w02b-p15-045-information}.
\item If \(\mu_1\ll\mu_2\) but \(\mu_2\) is not absolutely continuous
with respect to \(\mu_1\), then \(0<I(1:2)\leq\infty\) and
\(I(2:1)=J(1,2)=+\infty\). If neither measure is absolutely continuous
with respect to the other, \(I(1:2)=I(2:1)=+\infty\).
\end{enumerate}
\end{hblemma}

Part (2) shows that the extension changes nothing where the source's
definition applies. Clause (1) of the definition is the integral of the
log density ratio whenever that ratio exists as a Radon--Nikodym
derivative. Clause (2) assigns \(+\infty\) exactly when one observation
can be conclusive in favour of \(H_1\): an event of positive probability
under \(\mu_1\) has probability zero under \(\mu_2\). No inequality,
reduction or equality statement of
\cref{sec:p15-w02b-p15-045-inequalities,sec:p15-w02b-p15-045-sufficiency}
is asserted for the extended values; those statements keep the standing
hypotheses.

\section{Information inequalities and additivity}
\label{sec:p15-w02b-p15-045-inequalities}

\begin{hblemma}{thm:p15-w02b-p15-045-nonnegative}{proved=annex:p15-w02b-p15-045-nonnegative}{Lemma 3.1: almost positive definiteness}
Under the standing hypotheses of \cref{def:p15-w02b-p15-045-pair},
\(I(1:2)\geq0\), with equality if and only if \(f_1(x)=f_2(x)\)
\([\lambda]\), that is, if and only if \(\mu_1=\mu_2\). The same holds
for \(I(2:1)\), and \(J(1,2)\geq0\) with equality if and only if
\(\mu_1=\mu_2\).
\end{hblemma}

This is Lemma~3.1 of \cite[p.~82]{KullbackLeibler1951Information}. Its
proof writes \(I(1:2)=\int g\log g\,d\mu_2\) with \(g=f_1/f_2\), uses
\(\int g\,d\mu_2=1\), and controls \(\varphi(t)=t\log t\) by its
second-order expansion at \(t=1\). The annex follows that route and
shows that it remains valid when \(I(1:2)=+\infty\).

\begin{hblemma}{thm:p15-w02b-p15-045-event-bound}{proved=annex:p15-w02b-p15-045-event-bound}{Lemma 3.2: information from an event}
Under the standing hypotheses of \cref{def:p15-w02b-p15-045-pair}, for
every \(E\in\mathbf S\) with \(\lambda(E)>0\),
\[
  I_{1:2}(E)\geq\log\frac{\mu_1(E)}{\mu_2(E)},
\]
where both probabilities are positive. If \(I_{1:2}(E)<\infty\),
equality holds if and only if \(f_1(x)/f_2(x)=\mu_1(E)/\mu_2(E)\) for
\(\lambda\)-almost every \(x\in E\). If \(I_{1:2}(E)=\infty\) the
inequality is strict.
\end{hblemma}

The statement and the reduction to Lemma~3.1 through the normalized
measures \(\mu_i(\cdot\cap E)/\mu_i(E)\) are in
\cite[p.~82]{KullbackLeibler1951Information}. The source writes the
condition as \(\lambda(E)>0\); under the standing hypotheses this is
equivalent to \(\mu_1(E)>0\) and to \(\mu_2(E)>0\). The equality case is
added here.

\begin{hbdefinition}{def:p15-w02b-p15-045-product}{Independent components}
For \(k=1,2\) let \((X_k,\mathbf S_k)\) carry measures
\(\mu_1^{(k)},\mu_2^{(k)},\lambda_k\) and densities
\(f_1^{(k)},f_2^{(k)}\) satisfying the standing hypotheses of
\cref{def:p15-w02b-p15-045-pair}. On \(X=X_1\times X_2\) with
\(\mathbf S=\mathbf S_1\otimes\mathbf S_2\) put
\(\mu_i=\mu_i^{(1)}\otimes\mu_i^{(2)}\) and
\(\lambda=\lambda_1\otimes\lambda_2\)
(\cref{thm:p15-product-existence}). The coordinates \(x\in X_1\) and
\(y\in X_2\) are then independent under each \(\mu_i\). Write
\(I_x(1:2)\) and \(I_y(1:2)\) for the mean informations of the two
factor pairs and \(I_{xy}(1:2)\) for that of \((\mu_1,\mu_2)\).
\end{hbdefinition}

\begin{hbtheorem}{thm:p15-w02b-p15-045-additivity}{proved=annex:p15-w02b-p15-045-additivity}{Theorem 3.1: additivity for independent components}
In the setting of \cref{def:p15-w02b-p15-045-product}, the triple
\((\mu_1,\mu_2,\lambda)\) satisfies the standing hypotheses with
densities \(f_i(x,y)=f_i^{(1)}(x)f_i^{(2)}(y)\), and
\[
  I_{xy}(1:2)=I_x(1:2)+I_y(1:2)\qquad\text{in }[0,\infty].
\]
The same holds for \(I(2:1)\) and for \(J(1,2)\). In particular, for
\(n\) independent observations each distributed as \(\mu_i\) under
\(H_i\), the mean information of the \(n\)-tuple is \(n\,I(1:2)\).
\end{hbtheorem}

Theorem~3.1 and its proof (3.6) are in
\cite[pp.~82--83]{KullbackLeibler1951Information}. The source's proof
splits the integral of a sum into a sum of integrals without comment.
That step is legitimate here because each of the two terms has an
integrable negative part (\cref{thm:p15-w02b-p15-045-well-defined}),
so the identity holds even when one side is infinite.

\section{Statistics, sufficiency and data reduction}
\label{sec:p15-w02b-p15-045-sufficiency}

\begin{hbdefinition}{def:p15-w02b-p15-045-statistic}{Statistic, induced measures and sufficiency}
Under the standing hypotheses of \cref{def:p15-w02b-p15-045-pair},
let \((Y,\mathbf T)\) be a measurable space and
\(T:X\to Y\) a measurable map. Put
\[
  \nu_i(G)=\mu_i(T^{-1}G),\qquad
  \gamma(G)=\lambda(T^{-1}G)\qquad(G\in\mathbf T),
\]
and let \(T^{-1}(\mathbf T)=\{T^{-1}G:G\in\mathbf T\}\), a
sub-\(\sigma\)-algebra of \(\mathbf S\). Let \(g_i\) be a version of
\(d\nu_i/d\gamma\). The reduced quantities are
\[
  I'_{1:2}(G)=\frac1{\nu_1(G)}\int_G g_1(y)\log\frac{g_1(y)}{g_2(y)}\,d\gamma(y)
  \qquad(\nu_1(G)>0),
\]
\(I'_{1:2}(G)=0\) when \(\nu_1(G)=0\),
\(I'(1:2)=I'_{1:2}(Y)\), and \(I'(2:1)\), \(J'_{12}(G)\),
\(J'(1,2)\) correspondingly. The statistic \(T\) is \emph{sufficient}
for \(\{\mu_1,\mu_2\}\) if for every \(E\in\mathbf S\) there is a
\(T^{-1}(\mathbf T)\)-measurable function \(\phi_E\) with
\[
  \mu_i(E\cap A)=\int_A\phi_E\,d\mu_i
  \qquad(A\in T^{-1}(\mathbf T),\ i=1,2),
\]
that is, if one function is a version of the conditional probability of
\(E\) given \(T\) under both hypotheses.
\end{hbdefinition}

The induced measures and the reduced quantities are (4.1)--(4.3) of
\cite[p.~83]{KullbackLeibler1951Information}; there \(T\) is assumed to
map \(X\) onto \(Y\), which the arguments below never use. The source
takes the notion of sufficiency from Halmos and Savage
\cite[pp.~79 and~84]{KullbackLeibler1951Information}. The definition
above is the common-conditional-probability form, written out so that
the proofs below are self-contained. The finite binary case, where the
likelihood ratio is a minimal sufficient statistic, is
\cref{thm:p15-r43-sufficiency}.

\begin{hblemma}{thm:p15-w02b-p15-045-induced}{proved=annex:p15-w02b-p15-045-induced}{The induced pair and the conditional-expectation form of its densities}
In the setting of \cref{def:p15-w02b-p15-045-statistic}:
\begin{enumerate}
\item \(\nu_1\equiv\nu_2\), \(\gamma\) is a probability measure with
\(\gamma\equiv\{\nu_1,\nu_2\}\), and \(g_i\) may be chosen with
\(0<g_i(y)<\infty\) for every \(y\); with this choice
\((\nu_1,\nu_2,\gamma,g_1,g_2)\) satisfies the standing hypotheses on
\((Y,\mathbf T)\);
\item \(g_i\circ T\) is a version of the conditional expectation of
\(f_i\) given \(T^{-1}(\mathbf T)\) under \(\lambda\)
(\cref{def:p15-general-conditional-expectation}):
\(\int_{T^{-1}G}g_i(T(x))\,d\lambda(x)=\mu_i(T^{-1}G)\) for every
\(G\in\mathbf T\);
\item for every \(G\in\mathbf T\) with \(\nu_1(G)>0\),
\[
  I'_{1:2}(G)=\frac1{\mu_1(T^{-1}G)}\int_{T^{-1}G}d\mu_1(x)\log\frac{g_1T(x)}{g_2T(x)},
\]
\[
  I'(1:2)=\int d\mu_1(x)\log\frac{g_1T(x)}{g_2T(x)},
\]
where \(g_iT(x)\) abbreviates \(g_i(T(x))\).
\end{enumerate}
\end{hblemma}

Part (3) is (4.4) and (4.9) of
\cite[pp.~83--84]{KullbackLeibler1951Information}, where it is obtained
from Lemma~3 of Halmos and Savage. Part (2) is the form of that lemma
used in the proof below.

\begin{hblemma}{thm:p15-w02b-p15-045-density-criterion}{proved=annex:p15-w02b-p15-045-density-criterion}{The density step: sufficiency as measurability of the ratio}
In the setting of \cref{def:p15-w02b-p15-045-statistic}, with
\(0<g_i<\infty\) everywhere, the following are equivalent:
\begin{enumerate}
\item \(T\) is sufficient for \(\{\mu_1,\mu_2\}\);
\item there is a \(T^{-1}(\mathbf T)\)-measurable function \(W\) with
\(f_1(x)/f_2(x)=W(x)\) \([\lambda]\);
\item \(\dfrac{f_1(x)}{f_2(x)}=\dfrac{g_1T(x)}{g_2T(x)}\) \([\lambda]\).
\end{enumerate}
No finiteness of \(I(1:2)\) is assumed.
\end{hblemma}

Condition (2) is the source's (4.8). The source obtains (1)\(\Leftrightarrow\)(2)
from Corollary~2 of Halmos and Savage, and (1)\(\Rightarrow\)(3) by
the proof of their Theorem~1 together with their Lemma~2
\cite[p.~84]{KullbackLeibler1951Information}. The annex proves all
three implications directly for two measures with the same null sets.

\begin{hbtheorem}{thm:p15-w02b-p15-045-reduction}{proved=annex:p15-w02b-p15-045-reduction}{Theorem 4.1: data reduction, with the exact finiteness condition}
Assume the standing hypotheses of \cref{def:p15-w02b-p15-045-pair} and
let \(T\) be a statistic as in \cref{def:p15-w02b-p15-045-statistic}.
\begin{enumerate}
\item \(I(1:2)\geq I'(1:2)\), both sides in \([0,\infty]\).
\item If \(T\) is sufficient for \(\{\mu_1,\mu_2\}\), then
\(I(1:2)=I'(1:2)\), whether or not this common value is finite.
\item If \(I(1:2)<\infty\), then \(I'(1:2)<\infty\), and
\(I(1:2)=I'(1:2)\) if and only if \(T\) is sufficient for
\(\{\mu_1,\mu_2\}\).
\item The finiteness hypothesis in (3) cannot be removed: there are
pairs and statistics with \(I(1:2)=I'(1:2)=\infty\) for which \(T\)
is not sufficient (\cref{thm:p15-w02b-p15-045-infinite-equality}).
\end{enumerate}
\end{hbtheorem}

Kullback and Leibler state the theorem as ``\(I(1:2)\geq I'(1:2)\), with
equality if and only if \(T\) is a sufficient statistic'' and dispose of
the case \(I(1:2)=\infty\) as trivial
\cite[p.~83]{KullbackLeibler1951Information}. The case is trivial only
for the inequality. Part (4) shows that the ``only if'' direction fails
when the information is infinite, so the theorem is stated here with the
finiteness condition that the equality characterization needs. The
finite-case argument is the source's: on the set where both are
defined, \(I(1:2)-I'(1:2)\) is written as \(\int g\log g\,d\mu_{12}\) for a
probability measure \(\mu_{12}\) and \(g=f_1g_2T/(f_2g_1T)\), equations
(4.5)--(4.7) of \cite[pp.~83--84]{KullbackLeibler1951Information}, and
Lemma~3.1 is applied.

\begin{hbcorollary}{thm:p15-w02b-p15-045-nonsingular}{proved=annex:p15-w02b-p15-045-nonsingular}{Corollary 4.1 and the divergence form}
Assume the hypotheses of \cref{thm:p15-w02b-p15-045-reduction}.
\begin{enumerate}
\item If \(T\) is non-singular, that is, \(T^{-1}(\mathbf T)=\mathbf S\),
then \(I(1:2)=I'(1:2)\).
\item \(J(1,2)\geq J'(1,2)\); if \(T\) is sufficient, equality holds;
if \(J(1,2)<\infty\), equality holds if and only if \(T\) is
sufficient.
\end{enumerate}
\end{hbcorollary}

Part (1) is Corollary~4.1 of \cite[p.~84]{KullbackLeibler1951Information}.
Part (2) is an instance of the source's statement that each result of
its Sections~3 and~4 has an identical counterpart for \(J(1,2)\)
\cite[p.~85]{KullbackLeibler1951Information}. The finiteness condition
is again needed: the example of
\cref{thm:p15-w02b-p15-045-infinite-equality} has
\(J(1,2)=J'(1,2)=\infty\) with \(T\) not sufficient.

\begin{hbtheorem}{thm:p15-w02b-p15-045-events}{proved=annex:p15-w02b-p15-045-events}{Theorem 4.2: reduction event by event}
Assume the hypotheses of \cref{thm:p15-w02b-p15-045-reduction}.
\begin{enumerate}
\item If \(I_{1:2}(T^{-1}G)=I'_{1:2}(G)\) for all \(G\in\mathbf T\),
then \(I(1:2)=I'(1:2)\).
\item If \(T\) is sufficient, then \(I_{1:2}(T^{-1}G)=I'_{1:2}(G)\)
for all \(G\in\mathbf T\).
\item If \(I(1:2)<\infty\), then \(I_{1:2}(T^{-1}G)=I'_{1:2}(G)\) for all
\(G\in\mathbf T\) if and only if \(I(1:2)=I'(1:2)\).
\end{enumerate}
Without \(I(1:2)<\infty\), \(I(1:2)=I'(1:2)\) does not imply the
equalities for all \(G\) (\cref{thm:p15-w02b-p15-045-infinite-equality}).
\end{hbtheorem}

Theorem~4.2 of \cite[p.~84]{KullbackLeibler1951Information} asserts
the equivalence in (3) without a finiteness condition.

Section~5 of the source asserts that each result of its Sections~3
and~4 has an identical counterpart for \(J(1,2)\), as a consequence of
definitions (2.5) and (2.6), and gives no proofs
\cite[p.~85]{KullbackLeibler1951Information}. The same page records that
Jeffreys had noted the symmetry, positive definiteness, additivity and
invariance under non-singular transformations of \(J\). The counterparts
of Lemma~3.1, Theorem~3.1 and Theorem~4.1 are the \(J\) clauses of
\cref{thm:p15-w02b-p15-045-nonnegative},
\cref{thm:p15-w02b-p15-045-additivity} and
\cref{thm:p15-w02b-p15-045-nonsingular}(2). The remaining three follow.

\begin{hbtheorem}{thm:p15-w02b-p15-045-j-counterparts}{proved=annex:p15-w02b-p15-045-j-counterparts}{Section 5: the divergence counterparts of Lemma 3.2, Corollary 4.1 and Theorem 4.2}
Assume the standing hypotheses of \cref{def:p15-w02b-p15-045-pair}.
\begin{enumerate}
\item \emph{(Lemma~3.2 for \(J\).)} For every \(E\in\mathbf S\) with
\(\lambda(E)>0\),
\[
  J_{12}(E)\geq\log\frac{\mu_1(E)}{\mu_2(E)}+\log\frac{\mu_2(E)}{\mu_1(E)}=0 .
\]
If \(J_{12}(E)<\infty\), equality holds if and only if
\(f_1(x)/f_2(x)=\mu_1(E)/\mu_2(E)\) for \(\lambda\)-almost every
\(x\in E\). If \(J_{12}(E)=\infty\) the inequality is strict.
\end{enumerate}
Let \(T\) be a statistic as in \cref{def:p15-w02b-p15-045-statistic}.
\begin{enumerate}\setcounter{enumi}{1}
\item \emph{(Corollary~4.1 for \(J\).)} If \(T\) is non-singular, that
is, \(T^{-1}(\mathbf T)=\mathbf S\), then \(J(1,2)=J'(1,2)\).
\item \emph{(Theorem~4.2 for \(J\).)}
\begin{enumerate}
\item If \(J_{12}(T^{-1}G)=J'_{12}(G)\) for all \(G\in\mathbf T\), then
\(J(1,2)=J'(1,2)\).
\item If \(T\) is sufficient, then \(J_{12}(T^{-1}G)=J'_{12}(G)\) for
all \(G\in\mathbf T\).
\item If \(J(1,2)<\infty\), the following are equivalent:
\(J_{12}(T^{-1}G)=J'_{12}(G)\) for all \(G\in\mathbf T\);
\(J(1,2)=J'(1,2)\); \(T\) is sufficient.
\end{enumerate}
\end{enumerate}
Without \(J(1,2)<\infty\), \(J(1,2)=J'(1,2)\) implies neither the
equalities for all \(G\) nor sufficiency
(\cref{thm:p15-w02b-p15-045-infinite-equality}).
\end{hbtheorem}

In part (1) the two event bounds cancel, so the counterpart of
Lemma~3.2 for \(J\) is a nonnegativity statement on each event that does
not involve the probabilities of the event; its equality condition is the
one of Lemma~3.2 for either direction. Part (3)(c) carries the same
finiteness condition as \cref{thm:p15-w02b-p15-045-events}(3), and for
the same reason.

\section{Parametric families: Fisher's information and the normal case}
\label{sec:p15-w02b-p15-045-parametric}

For two neighbouring members of a parametric family, Kullback and
Leibler state that \(I\) and \(J\) agree with a quadratic form in the
parameter increment whose coefficients are the elements of Fisher's
information matrix, to within second-order terms. They state this ``with
suitable assumptions on the density function'' and refer to p.~774 of
a 1934 paper of Doob for such assumptions, without listing them
\cite[p.~81, eqs.~(2.7)--(2.9)]{KullbackLeibler1951Information}. The
conditions (R1)--(R3) below are one explicit set of sufficient
conditions. They are supplied here, not taken from the source.

\begin{hbdefinition}{def:p15-w02b-p15-045-regular-family}{A regular parametric family and Fisher's information matrix}
Let \((X,\mathbf S)\) be a measurable space, \(\lambda\) a
\(\sigma\)-finite measure on \(\mathbf S\), \(k\geq1\), and
\(\Theta\subseteq\mathbb R^k\) open. For each \(\theta\in\Theta\) let
\(f(\cdot,\theta)\) be \(\mathbf S\)-measurable with
\(0<f(x,\theta)<\infty\) for every \(x\in X\) and
\(\int f(x,\theta)\,d\lambda(x)=1\), and let
\(\mu_\theta(E)=\int_Ef(x,\theta)\,d\lambda(x)\). Write
\(\partial_\alpha=\partial/\partial\theta_\alpha\), \(|\cdot|\) for the
Euclidean norm, and \(\ell(x,\theta)=\log f(x,\theta)\).
The family is \emph{regular at} \(\theta\in\Theta\) if there is \(r>0\)
such that the closed ball \(B=\{\theta':|\theta'-\theta|\leq r\}\) lies
in \(\Theta\) and:
\begin{enumerate}
\item[(R1)] for every \(x\in X\), \(\theta'\mapsto f(x,\theta')\) is
twice continuously differentiable on an open set containing \(B\);
\item[(R2)] there is a \(\lambda\)-integrable \(N\geq0\) such that
\[
  |\partial_\alpha f(x,\theta')|\leq N(x),\qquad
  |\partial_\alpha\partial_\beta f(x,\theta')|\leq N(x)
\]
for all \(x\in X\), \(\theta'\in B\) and \(\alpha,\beta=1,\dots,k\);
\item[(R3)] there is a \(\lambda\)-integrable \(K\geq0\) with
\(f(x,\theta')\,|\partial_\alpha\partial_\beta\ell(x,\theta'')|\leq K(x)\)
for all \(x\in X\), \(\theta',\theta''\in B\) and
\(\alpha,\beta=1,\dots,k\).
\end{enumerate}
\emph{Fisher's information matrix} at \(\theta\) has the elements
\[
  g_{\alpha\beta}=\int f\Bigl(\frac1f\frac{\partial f}{\partial\theta_\alpha}\Bigr)
                     \Bigl(\frac1f\frac{\partial f}{\partial\theta_\beta}\Bigr)d\lambda,
  \qquad\alpha,\beta=1,\dots,k,
\]
with \(f\) and its derivatives evaluated at \((x,\theta)\). For
\(\Delta\theta\in\mathbb R^k\) with \(\theta+\Delta\theta\in\Theta\),
\(I(\theta;\theta+\Delta\theta)\) and \(J(\theta,\theta+\Delta\theta)\)
denote \(I(1:2)\) and \(J(1,2)\) for \(f_1=f(\cdot,\theta)\) and
\(f_2=f(\cdot,\theta+\Delta\theta)\).
\end{hbdefinition}

The display is the source's (2.9), and \(I(\theta;\theta+\Delta\theta)\),
\(J(\theta,\theta+\Delta\theta)\) are its notation. Because every
density is positive, \(\mu_\theta(E)=0\) exactly when \(\lambda(E)=0\).
So any two members of the family satisfy the standing hypotheses with the
probability measure \(\mu_\theta\) as dominating measure, and by
\cref{thm:p15-w02b-p15-045-well-defined}(2) the quantities may be
computed with \(\lambda\).

\begin{hbtheorem}{thm:p15-w02b-p15-045-fisher}{proved=annex:p15-w02b-p15-045-fisher}{Equations (2.7)--(2.9): the second-order relation to Fisher's information}
Let the family of \cref{def:p15-w02b-p15-045-regular-family} be regular
at \(\theta\), with radius \(r\). All sums below run over
\(\alpha,\beta=1,\dots,k\).
\begin{enumerate}
\item Every \(g_{\alpha\beta}\) is finite, \(g_{\alpha\beta}=g_{\beta\alpha}\),
\(\sum g_{\alpha\beta}v_\alpha v_\beta\geq0\) for every \(v\in\mathbb R^k\), and
\[
  g_{\alpha\beta}=-\int f(x,\theta)\,\partial_\alpha\partial_\beta\ell(x,\theta)\,d\lambda(x).
\]
\item For \(0<|\Delta\theta|<r\), the quantities
\(I(\theta;\theta+\Delta\theta)\), \(I(\theta+\Delta\theta;\theta)\) and
\(J(\theta,\theta+\Delta\theta)\) are finite, and as
\(\Delta\theta\to0\),
\[
  \frac{I(\theta;\theta+\Delta\theta)-\tfrac12\sum g_{\alpha\beta}\Delta\theta_\alpha\Delta\theta_\beta}{|\Delta\theta|^2}\to0,
\]
\[
  \frac{J(\theta,\theta+\Delta\theta)-\sum g_{\alpha\beta}\Delta\theta_\alpha\Delta\theta_\beta}{|\Delta\theta|^2}\to0.
\]
\end{enumerate}
\end{hbtheorem}

Part (2) is the source's (2.7) and (2.8), with ``to within second order
terms'' made precise as a remainder that is small compared with
\(|\Delta\theta|^2\). (R2) is used to differentiate
\(\int f\,d\lambda=1\) twice under the integral sign. This gives the
zero mean of each \(\partial_\alpha\ell\) and the identity in (1). (R3)
dominates the Taylor remainder of \(\ell\) near \(\theta\), under
\(\mu_\theta\) and also under the shifted measure
\(\mu_{\theta+\Delta\theta}\), which the reverse direction needs. The
statement is local. It asserts nothing about \(I\) or \(J\) for
parameter values outside \(B\), and it does not make \(J\) a distance:
\cref{thm:p15-w02b-p15-045-normal} gives a family in which
\(I(\theta;\theta+\Delta\theta)\) is exactly a quadratic form in
\(\Delta\theta\), and \(J\) still violates the triangle inequality
(\cref{thm:p15-w02b-p15-045-not-metric}(2)).

\begin{hbtheorem}{thm:p15-w02b-p15-045-normal}{proved=annex:p15-w02b-p15-045-normal}{Equation (2.10): normal populations with a common covariance matrix}
Let \(k\geq1\), let \(A\) be an invertible real \(k\times k\) matrix,
put \(\Sigma=AA^{\mathsf T}\), and let \(\sigma^{\alpha\beta}\) be the
elements of \(\Sigma^{-1}\). Let \(\rho\) be the \(k\)-fold product of
the standard normal law (\cref{def:p15-normal-law}) on the Borel sets of
\(\mathbb R^k\), and for \(m_1,m_2\in\mathbb R^k\) let \(\mu_i\) be the
image of \(\rho\) under \(z\mapsto m_i+Az\); that is, \(\mu_i\) is the
law of \(m_i+AZ\) when \(Z\) has law \(\rho\). Put
\(\delta=m_1-m_2\), with components \(\delta_\alpha\). Then:
\begin{enumerate}
\item \(\mu_i\) has mean vector \(m_i\) and covariance matrix
\(\Sigma\);
\item \(\mu_1\equiv\mu_2\), so the standing hypotheses hold with
\(\lambda=\tfrac12(\mu_1+\mu_2)\);
\item \(I(1:2)=I(2:1)=\tfrac12\sum\delta_\alpha\delta_\beta\sigma^{\alpha\beta}\)
and
\[
  J(1,2)=\sum\delta_\alpha\delta_\beta\sigma^{\alpha\beta},
  \qquad\alpha,\beta=1,\dots,k;
\]
\item \(J(1,2)=|A^{-1}m_1-A^{-1}m_2|^2\).
\end{enumerate}
\end{hbtheorem}

Part (3) is (2.10) of \cite[p.~81]{KullbackLeibler1951Information},
stated there for multivariate normal populations with a common matrix of
variances and covariances, with \(\delta_\alpha\) the differences of the
population means. The source identifies this \(J(1,2)\) with \(k\) times
Mahalanobis' generalized distance \cite[pp.~81--82]{KullbackLeibler1951Information}.
The proof uses only results of this chapter. The affine map
\(x\mapsto A^{-1}(x-m_2)\) is non-singular, so by Corollary~4.1
(\cref{thm:p15-w02b-p15-045-nonsingular}) it preserves information. It
carries the pair to two products of one-dimensional normal laws with
unit variance. Additivity (\cref{thm:p15-w02b-p15-045-additivity}) and
the one-dimensional computation of
\cref{thm:p15-w02b-p15-045-not-metric}(2) then finish the calculation.
By part (4), \(\sqrt{J(1,2)}\) is a Euclidean distance between
transformed mean vectors when \(\Sigma\) is held fixed. \(J(1,2)\)
itself is the square of that distance and is not a metric. With
\(\mu_\theta\) the law of \(\theta+AZ\), this is a family in which
\(\tfrac12\sum\sigma^{\alpha\beta}\Delta\theta_\alpha\Delta\theta_\beta\)
equals \(I(\theta;\theta+\Delta\theta)\) exactly, for every increment.

\section{What relative information is not}
\label{sec:p15-w02b-p15-045-distinctions}

The word ``information'' and the word ``divergence'' invite three
misreadings: that \(I(1:2)\) or \(J(1,2)\) is a distance between the
two measures, that \(\log(f_1/f_2)\) is a posterior probability, and
that \(I(1:2)\) is the expected utility of observing. Each is false
under the standing hypotheses, and the next three results make the
differences exact.

\begin{hbtheorem}{thm:p15-w02b-p15-045-bayes}{proved=annex:p15-w02b-p15-045-bayes}{The Bayes reading of the information in an observation}
Assume the standing hypotheses of \cref{def:p15-w02b-p15-045-pair}.
Let \(\alpha_1,\alpha_2>0\) with \(\alpha_1+\alpha_2=1\) be prior
probabilities of \(H_1,H_2\), so that the hypothesis is drawn with these
probabilities and then \(x\) is drawn from \(\mu_i\). Then
\[
  P(H_i\mid x)=\frac{\alpha_if_i(x)}{\alpha_1f_1(x)+\alpha_2f_2(x)}
\]
is a version of the posterior probability of \(H_i\) given \(x\), each
\(P(H_i\mid x)\) lies in \((0,1)\) \([\lambda]\), and
\[
  \log\frac{f_1(x)}{f_2(x)}
   =\log\frac{P(H_1\mid x)}{P(H_2\mid x)}-\log\frac{\alpha_1}{\alpha_2}
   \qquad[\lambda].
\]
Consequently \(I(1:2)\) is the mean, under \(H_1\), of the change in log
odds produced by one observation. It does not depend on
\(\alpha_1,\alpha_2\), and it can be \(+\infty\) although every
posterior probability lies strictly between \(0\) and \(1\)
(\cref{thm:p15-w02b-p15-045-footnote}).
\end{hbtheorem}

The identity is footnote~2 of
\cite[p.~80]{KullbackLeibler1951Information}. It shows what kind of
quantity the information in \(x\) is: a log Bayes factor, the additive
increment that turns prior log odds into posterior log odds. It is not a
probability of either hypothesis. A large \(I(1:2)\) is compatible with
a small posterior probability of \(H_1\) when \(\alpha_1\) is small,
because \(I(1:2)\) does not involve the prior.

\begin{hbtheorem}{thm:p15-w02b-p15-045-not-metric}{proved=annex:p15-w02b-p15-045-not-metric}{Neither \(I(1:2)\) nor \(J(1,2)\) is a metric}
\begin{enumerate}
\item On \(X=\{0,1\}\) with \(\lambda\) counting measure, let \(\mu_1\)
give mass \(\tfrac12,\tfrac12\) and \(\mu_2\) mass
\(\tfrac14,\tfrac34\) to \(0,1\). Then
\(I(1:2)=\tfrac12\log\tfrac43\approx0.1438\) and
\(I(2:1)=\tfrac14\log\tfrac{27}{16}\approx0.1308\), so \(I\) is not
symmetric.
\item On \(\mathbb R\) with Lebesgue measure as dominating measure, let
\(\mu_a\) be the normal law with mean \(a\) and variance \(1\). For the
pair \((\mu_a,\mu_b)\), \(I(1:2)=I(2:1)=\tfrac12(a-b)^2\) and
\(J(1,2)=(a-b)^2\). For \(a=0\), \(b=1\), \(c=2\) the divergence of the
pair \((\mu_0,\mu_2)\) is \(4\), which exceeds the sum \(2\) of the
divergences of \((\mu_0,\mu_1)\) and \((\mu_1,\mu_2)\). So \(J\) is
symmetric but violates the triangle inequality, and so does \(I\).
\end{enumerate}
\end{hbtheorem}

Part (2) is the case \(k=1\) of the source's formula (2.10)
\cite[p.~81]{KullbackLeibler1951Information}, computed directly; the
general case is \cref{thm:p15-w02b-p15-045-normal}. The
source calls \(J(1,2)\) the divergence, borrowing Jeffreys' term, and
places ``distance'' in quotation marks
\cite[pp.~79 and~81]{KullbackLeibler1951Information}. Its Section~5
notes that \(J(1,2)\) is symmetric in \(\mu_1,\mu_2\) and independent of
the prior probabilities \cite[p.~85]{KullbackLeibler1951Information}. The
finite version of the non-metric remark is in \cref{def:p16-rat42-kl}.

\begin{hbwarning}{def:p15-w02b-p15-045-not-distance}{Divergence is not a distance}
Under the standing hypotheses, \(I(1:2)\) is directed: it averages over
\(H_1\) and can differ from \(I(2:1)\), even to the extent that one is
finite and the other infinite (\cref{thm:p15-w02b-p15-045-footnote}).
\(J(1,2)\) is symmetric and vanishes only at \(\mu_1=\mu_2\), but it
violates the triangle inequality (\cref{thm:p15-w02b-p15-045-not-metric}).
Neither quantity may be used as a metric. Neither is defined by
\cref{def:p15-w02b-p15-045-information} for a pair whose null sets
differ; for such a pair the values of
\cref{def:p15-w02b-p15-045-extension} are extensions, and at least one
direction is \(+\infty\) (\cref{thm:p15-w02b-p15-045-extension}(3)).
\end{hbwarning}

\begin{hbtheorem}{thm:p15-w02b-p15-045-not-utility}{proved=annex:p15-w02b-p15-045-not-utility}{More information can have lower expected utility}
Let \(X_1=X_2=\{0,1\}\) with counting measures, and take the product
setting of \cref{def:p15-w02b-p15-045-product} with
\[
  \mu_1^{(1)}=\bigl(\tfrac12,\tfrac12\bigr),\quad
  \mu_2^{(1)}=\bigl(\tfrac{999}{1000},\tfrac1{1000}\bigr),\qquad
  \mu_1^{(2)}=\bigl(\tfrac9{10},\tfrac1{10}\bigr),\quad
  \mu_2^{(2)}=\bigl(\tfrac1{10},\tfrac9{10}\bigr),
\]
listing the masses of \(0\) and \(1\). Let \(T(x,y)=x\) and
\(T'(x,y)=y\).
\begin{enumerate}
\item The reduced informations are
\(I_x(1:2)=\tfrac12\log\tfrac{250000}{999}\approx2.761\) and
\(I_y(1:2)=\tfrac45\log9\approx1.758\).
\item With prior probabilities \(\tfrac12,\tfrac12\), acts ``assert
\(H_1\)'' and ``assert \(H_2\)'', and utility \(1\) for a true assertion
and \(0\) otherwise, the largest expected utility of a rule based on
\(T\) is \(\tfrac{1499}{2000}=0.7495\), and of a rule based on
\(T'\) is \(\tfrac9{10}\).
\end{enumerate}
So the statistic carrying more information has the lower expected
utility in this decision problem.
\end{hbtheorem}

\(I(1:2)\) is an expectation, under \(H_1\) alone, of a function fixed by
the two hypotheses. A decision problem adds what \(I(1:2)\) lacks: a set
of acts, a utility on acts and states, and prior probabilities over the
hypotheses. An expected utility is an average over both hypotheses with
those prior weights. Ranking experiments or statistics by \(I(1:2)\) is
therefore a different criterion from ranking them by their expected
utility for a given decision problem, and the two rankings can
disagree. The source's own applications in its Section~6 choose a
critical region (\cref{thm:p15-w02b-p15-045-critical}) and a grouping
(\cref{def:p15-w02b-p15-045-grouping}) by the information criterion
\cite[pp.~85--86]{KullbackLeibler1951Information}; that choice is a
separate decision about which criterion to use, not a consequence of
decision theory.

\section{Worked cases}
\label{sec:p15-w02b-p15-045-cases}

\begin{hbtheorem}{thm:p15-w02b-p15-045-footnote}{proved=annex:p15-w02b-p15-045-footnote}{Positive densities, infinite information}
Let \(X=(0,1)\) with its Borel sets, \(\mu_1=\lambda\) Lebesgue measure on
\((0,1)\), and \(f_1=1\). Let
\[
  k^{-1}=\int_0^1e^{-1/t}\,dt,\qquad f_2(x)=ke^{-1/x}\quad(0<x<1),
\]
and \(\mu_2(E)=\int_Ef_2\,d\lambda\). Then the standing hypotheses hold,
with \(0<f_i(x)<\infty\) for every \(x\), and
\begin{enumerate}
\item \(I(1:2)=+\infty\);
\item \(0\leq I(2:1)\leq\log k<\infty\); numerical quadrature gives
\(\log k\approx1.907\) and \(I(2:1)\approx0.430\);
\item for any prior probabilities \(\alpha_1,\alpha_2>0\), every
posterior probability \(P(H_i\mid x)\) of
\cref{thm:p15-w02b-p15-045-bayes} lies in \((0,1)\) for every
\(x\in(0,1)\).
\end{enumerate}
\end{hbtheorem}

The pair is the source's footnote~3 example, which shows that the
standing hypotheses do not imply finite information
\cite[p.~80, n.~3]{KullbackLeibler1951Information}. The source states
only that \(I(1:2)\) is infinite. Parts (2) and (3) add that the reverse
information is finite and that no observation is conclusive. The
infinite value therefore comes from the unbounded tail of the log ratio
\(1/x-\log k\) near \(0\), not from a support mismatch. A countable
analogue, with both directions computed exactly, is
\cref{thm:p15-infinite-entropy-finite-scores}.

\begin{hbtheorem}{thm:p15-w02b-p15-045-support}{proved=annex:p15-w02b-p15-045-support}{Support mismatch: the extended values as limits}
Let \(\mu_1,\mu_2\) be probability measures on \((X,\mathbf S)\), and
suppose \(\mu_1(E)>0=\mu_2(E)\) for some \(E\in\mathbf S\). So the
standing hypotheses fail, and \(I(1:2)=+\infty\) by
\cref{def:p15-w02b-p15-045-extension}(2).
\begin{enumerate}
\item \emph{No assumption on the sets that \(\mu_2\) charges.} For
\(0<\varepsilon<1\) put
\(\bar\mu_1^\varepsilon=(1-\varepsilon)\mu_1+\varepsilon\mu_2\) and
\(\bar\mu_2^\varepsilon=(1-\varepsilon)\mu_2+\varepsilon\mu_1\). Then
\(\bar\mu_1^\varepsilon\equiv\bar\mu_2^\varepsilon\), the pair satisfies
the standing hypotheses with \(\lambda=\tfrac12(\mu_1+\mu_2)\), and its
information \(\bar I^\varepsilon(1:2)\) satisfies
\[
  \bar I^\varepsilon(1:2)\geq(1-\varepsilon)\,\mu_1(E)\log\frac{1-\varepsilon}\varepsilon-\frac1e .
\]
Hence \(\bar I^\varepsilon(1:2)\to\infty\) as \(\varepsilon\to0\).
\item \emph{The one-sided case.} Suppose in addition that
\(\mu_2\ll\mu_1\), and for \(0<\varepsilon<1\) put
\(\mu_2^\varepsilon=(1-\varepsilon)\mu_2+\varepsilon\mu_1\). Then
\(\mu_1\equiv\mu_2^\varepsilon\), the pair
\((\mu_1,\mu_2^\varepsilon)\) satisfies the standing hypotheses with
\(\lambda=\mu_1\), and its informations \(I^\varepsilon(1:2)\) and
\(I^\varepsilon(2:1)\) satisfy
\[
  I^\varepsilon(1:2)\geq\mu_1(E)\log\frac1\varepsilon+\mu_1(X\setminus E)\log\mu_1(X\setminus E),
\]
the last term read as \(0\) when \(\mu_1(X\setminus E)=0\), so
\(I^\varepsilon(1:2)\to\infty\); and
\[
  I^\varepsilon(2:1)\leq(1-\varepsilon)\,I(2:1),\qquad
  I^\varepsilon(2:1)\to I(2:1)\quad(\varepsilon\to0),
\]
where \(I(2:1)=\int\log(d\mu_2/d\mu_1)\,d\mu_2\in(0,\infty]\) is the
value of \cref{def:p15-w02b-p15-045-extension}(1) with the indices
exchanged.
\item For \(X=(0,2)\), \(\mu_1\) uniform on \((0,2)\) and \(\mu_2\)
uniform on \((0,1)\), \(I(2:1)=\log2\) and, in the notation of (2),
\[
  I^\varepsilon(1:2)=-\tfrac12\log\bigl(\varepsilon(2-\varepsilon)\bigr)\to\infty,
\]
\[
  I^\varepsilon(2:1)=\bigl(1-\tfrac\varepsilon2\bigr)\log(2-\varepsilon)
     +\tfrac\varepsilon2\log\varepsilon\to\log2 .
\]
\end{enumerate}
\end{hbtheorem}

Part (1) is the justification for
\cref{def:p15-w02b-p15-045-extension}(2) for every pair in which
\(\mu_1\) charges an event that \(\mu_2\) does not, whatever the events
charged by \(\mu_2\) alone. Mixing each measure with a small multiple of
the other gives a pair with the same null sets and with
\(|\bar\mu_i^\varepsilon(A)-\mu_i(A)|\leq\varepsilon\) for every
\(A\in\mathbf S\), and its information tends to \(+\infty\). Part (2) justifies
\cref{def:p15-w02b-p15-045-extension}(1) in the direction in which
absolute continuity holds. The extended value \(I(2:1)\) is the limit of
the informations of mutually absolutely continuous approximations, and
each approximation has at most that information. A support mismatch with
\(\mu_2\ll\mu_1\) is thus an extreme asymmetry: \(I(1:2)=+\infty\) while
\(I(2:1)\) may be finite, as in (3), where it is \(\log2\). The lower
bounds in (1) and (2) are Lemma~3.2 applied to \(E\) and to its
complement.

\begin{hbtheorem}{thm:p15-w02b-p15-045-loss}{proved=annex:p15-w02b-p15-045-loss}{Information loss by a statistic that is not sufficient}
In the setting of \cref{thm:p15-w02b-p15-045-not-utility}, the
statistic \(T(x,y)=x\) is not sufficient for \(\{\mu_1,\mu_2\}\), and
\[
  I'(1:2)=I_x(1:2)\approx2.761<I(1:2)=I_x(1:2)+I_y(1:2)\approx4.519 .
\]
\end{hbtheorem}

\begin{hbtheorem}{thm:p15-w02b-p15-045-infinite-equality}{proved=annex:p15-w02b-p15-045-infinite-equality}{Equal infinite information without sufficiency}
Take the product setting of \cref{def:p15-w02b-p15-045-product} with
the pair of \cref{thm:p15-w02b-p15-045-footnote} on \(X_1=(0,1)\) and the
pair \(\mu_1^{(2)}=(\tfrac9{10},\tfrac1{10})\),
\(\mu_2^{(2)}=(\tfrac1{10},\tfrac9{10})\) on \(X_2=\{0,1\}\), and let
\(T(x,y)=x\).
\begin{enumerate}
\item \(I(1:2)=I'(1:2)=+\infty\) and \(J(1,2)=J'(1,2)=+\infty\), but
\(T\) is not sufficient.
\item For \(G=(\tfrac12,1)\),
\(I_{1:2}(T^{-1}G)=I'_{1:2}(G)+\tfrac45\log9\neq I'_{1:2}(G)\), and
\(I'_{1:2}(G)\) is finite.
\item For the same \(G\),
\(J_{12}(T^{-1}G)=J'_{12}(G)+\tfrac85\log9\neq J'_{12}(G)\), and
\(J'_{12}(G)\) is finite.
\end{enumerate}
\end{hbtheorem}

This example is what limits the equality statements in
\cref{thm:p15-w02b-p15-045-reduction,thm:p15-w02b-p15-045-nonsingular,thm:p15-w02b-p15-045-events,thm:p15-w02b-p15-045-j-counterparts}
to finite information. It is built from the source's own footnote~3 pair.

\begin{hbtheorem}{thm:p15-w02b-p15-045-critical}{proved=annex:p15-w02b-p15-045-critical}{Section 6(1): comparing two critical regions}
Let \(0<\alpha<1\), \(0<c<1\) and \(0<p<1\). Let \(X\) be a countable
set containing \(-2\), \(0\) and \(2\), with \(\mathbf S\) the set of
all subsets, and let \(\mu_1,\mu_2\) be probability measures on
\(\mathbf S\) that give positive mass to every point of \(X\) and satisfy
\[
  \mu_1(\{-2\})=\mu_1(\{2\})=\frac\alpha2,\qquad
  \mu_2(\{-2\})=pc,\qquad
  \mu_2(\{2\})=(1-p)c,
\]
\[
  \frac{\mu_1(\{0\})}{\mu_2(\{0\})}=\frac{1-\alpha}{1-c}.
\]
Let \(F=\{-2,2\}\) and \(G=\{0\}\).
\begin{enumerate}
\item The standing hypotheses hold, and
\[
  I_{1:2}(F)=\frac1\alpha\Bigl(\frac\alpha2\log\frac\alpha{2pc}
             +\frac\alpha2\log\frac\alpha{2c(1-p)}\Bigr),\qquad
  I_{1:2}(G)=\log\frac{1-\alpha}{1-c}.
\]
\item
\(\displaystyle I_{1:2}(F)-I_{1:2}(G)=\log\frac{\alpha(1-c)}{2c(1-\alpha)}-\frac12\log\bigl(p(1-p)\bigr)\).
\item \(I_{1:2}(G)<I_{1:2}(F)\) if and only if
\(\alpha(1-c)>2c(1-\alpha)\sqrt{p(1-p)}\). In particular the inequality
holds whenever \(c<\alpha\), and whenever \(c=\alpha\) and
\(p\neq\tfrac12\); it fails whenever \(p=\tfrac12\) and \(c\geq\alpha\).
\end{enumerate}
\end{hbtheorem}

Section~6(1) of \cite[p.~85]{KullbackLeibler1951Information} applies
the event information \(I_{1:2}(E)\) to a testing problem of Lehmann,
which the source cites at p.~2 of Lehmann's 1950 paper on principles of
testing hypotheses \cite[pp.~85--86]{KullbackLeibler1951Information}.
Index \(1\) refers to Lehmann's hypothesis \(H\), index \(2\) to any of
the alternatives, and the candidate critical regions are \(F=\{-2,2\}\)
and \(G=\{0\}\). The source displays the two formulas of part (1) as its
(6.1), with the second written as \(\frac1\alpha\cdot\alpha\log\frac{1-\alpha}{1-c}\).
It says that \(I_{1:2}(G)<I_{1:2}(F)\) can be verified readily, and
concludes that \(G\) should be the critical region. The source does not
reproduce Lehmann's distributions or the admissible values of
\(\alpha\), \(c\) and \(p\). The masses in the hypotheses above are those
that (6.1) presupposes when each line is read term by term through
definition (2.3): the first line through the two atoms of \(F\), and the
second through the single atom of \(G\), for which (2.3) reduces to
\(\log(\mu_1(\{0\})/\mu_2(\{0\}))\). Part (3) shows that the inequality
is not automatic. It depends on \((\alpha,c,p)\), and it holds for every
\(p\) once \(c<\alpha\). The selection rule itself is the source's:
reject \(H\) on the region where an observation carries the smaller
mean information in favour of \(H_1\). It is not derived from the error
probabilities of the two tests (\cref{sec:p15-w02b-p15-045-distinctions}).

\begin{hbexample}{def:p15-w02b-p15-045-grouping}{The source's grouped multinomial test}
Kullback and Leibler compare a multinomial population
\((p_1,\dots,p_k)\), \(p_1\geq\dots\geq p_k\), with the uniform population
\((1/k,\dots,1/k)\) when the observations must be grouped into two
categories \cite[pp.~85--86]{KullbackLeibler1951Information}. The
source does not say whether some \(p_i\) may vanish. The comparison
below assumes in addition that \(p_k>0\), so that every \(p_i\) is
positive. Then on \(X=\{1,\dots,k\}\) with all subsets, \(\mu_1\) giving
mass \(p_i\) to \(i\), \(\mu_2\) the uniform probability and
\(\lambda=\mu_2\), the standing hypotheses of
\cref{def:p15-w02b-p15-045-pair} hold with \(f_1(i)=kp_i\) and
\(f_2(i)=1\), and \(J\) below is a finite sum. Grouping
the first \(r\) categories against the rest is the statistic that maps
category \(i\) to \(1\) if \(i\leq r\) and to \(2\) otherwise. Its reduced
divergence is
\[
  J'=\Bigl(\sum_{i=1}^rp_i-\frac rk\Bigr)\log\frac{\sum_{i=1}^rp_i}{r/k}
    +\Bigl(\sum_{i=r+1}^kp_i-\frac{k-r}k\Bigr)\log\frac{\sum_{i=r+1}^kp_i}{(k-r)/k},
\]
where for \(r=k\) the second class is empty, the statistic is constant,
the second term is absent and \(J'=0\). The full divergence is
\(J=\sum_{i=1}^k(p_i-1/k)\log(kp_i)\). By
\cref{thm:p15-w02b-p15-045-nonsingular}(2), \(J'\leq J\). Because
\(J<\infty\), the same result gives equality if and only if the grouping
is sufficient. By \cref{thm:p15-w02b-p15-045-density-criterion} that
happens exactly when \(f_1/f_2=kp_i\) is a function of the class of
\(i\), that is, when \(p_1=\dots=p_r\) and \(p_{r+1}=\dots=p_k\). By
additivity (\cref{thm:p15-w02b-p15-045-additivity}), \(n\) grouped
observations carry as much divergence as \(N\) ungrouped ones when
\(nJ'=NJ\), so, for \(p\) not uniform (so that \(0<J<\infty\)),
\(J'/J\) measures the efficiency of the grouped test.

The positivity assumption is not idle. If \(p_k=0\), the uniform
population charges category \(k\) and \(\mu_1\) does not, so the
standing hypotheses fail and \cref{def:p15-w02b-p15-045-information}
defines neither \(J\) nor \(J'\). The extended values of
\cref{def:p15-w02b-p15-045-extension} do not rescue the equality
statement. Take \(k=3\), \(p=(0.6,0.4,0)\) and \(r=2\). Since
\(\mu_2(\{3\})=\tfrac13>0=\mu_1(\{3\})\), the extended \(I(2:1)\) is
\(+\infty\), and so is \(J\). The grouped pair gives mass \((1,0)\) and
\((\tfrac23,\tfrac13)\) to the two classes, so the extended \(I'(2:1)\)
and \(J'\) are \(+\infty\) as well. Thus \(J'=J\), although
\(kp_i\) takes the two values \(1.8\) and \(1.2\) on the class
\(\{1,2\}\), so the ratio of the two populations is not a function of the
grouping. Equality of extended values therefore says nothing about
whether a grouping loses information, which is why
\cref{thm:p15-w02b-p15-045-nonsingular} keeps the standing hypotheses
and the condition \(J<\infty\).

For \(p=(.5,.3,.1,.1)\) and base-\(10\) logarithms the source prints
\(J=0.1986\) and, for \(r=1,2,3,4\), the values \(.1193\), \(.0903\),
\(.0716\), \(0.0\); it concludes that the best grouping is \(r=1\),
with efficiency \(.1193/.1986=.6007\), a loss of \(40\) percent.
Evaluating the displayed formula gives \(J=0.19860\) and
\(J'=0.11928\), \(0.18062\), \(0.07157\), \(0\) for \(r=1,2,3,4\).
For \(r=2\), \(\sum_{i\leq2}p_i=0.8\), and
\[
  J'=0.3\log_{10}1.6+0.3\log_{10}2.5=0.3\log_{10}4=0.18062,
\]
twice the printed value. With the formula as displayed, the best
two-category grouping is therefore the one that pools \(p_1\) with
\(p_2\) and \(p_3\) with \(p_4\). Its efficiency is
\(0.18062/0.19860=0.9095\), a loss of about \(9\) percent.
Every grouping loses information here: \(J'<J\) for each \(r\), as
\cref{thm:p15-w02b-p15-045-nonsingular} requires, because for no \(r\)
is the ratio \(p_i/(1/k)\) constant on each class of the grouping.
\end{hbexample}

\section{Procedure}
\label{sec:p15-w02b-p15-045-procedure}

\begin{hbmethod}{def:p15-w02b-p15-045-procedure}{Computing and comparing relative information}
\begin{enumerate}
\item \emph{Declare the pair.} State the sample space, its
\(\sigma\)-algebra, and the two probability measures, including which
one is \(\mu_1\); the quantity \(I(1:2)\) is directed.
\item \emph{Check null sets.} Verify \(\mu_1\equiv\mu_2\). If the null
sets differ, \cref{def:p15-w02b-p15-045-information} does not apply, and
only the extended values of \cref{def:p15-w02b-p15-045-extension} are
available, reported as extensions. If an event has positive probability
under \(\mu_1\) and probability zero under \(\mu_2\), whatever \(\mu_2\)
charges, the extended \(I(1:2)\) is \(+\infty\)
(\cref{thm:p15-w02b-p15-045-support}(1)), and one such observation is
conclusive. If moreover \(\mu_2\ll\mu_1\), compute the finite or
infinite \(I(2:1)=\int\log(d\mu_2/d\mu_1)\,d\mu_2\)
(\cref{thm:p15-w02b-p15-045-support}(2)). The results of
\cref{sec:p15-w02b-p15-045-inequalities,sec:p15-w02b-p15-045-sufficiency}
are not applied to extended values.
\item \emph{Choose a dominating measure.} Any \(\sigma\)-finite measure
with the same null sets will do
(\cref{thm:p15-w02b-p15-045-well-defined}); compute the densities and
the ratio \(f_1/f_2\).
\item \emph{Integrate with the sign convention.} Compute
\(\int f_1\log(f_1/f_2)\,d\lambda\) as an extended integral; its
negative part is always finite, and the value \(+\infty\) must be
reported as such, not discarded.
\item \emph{For independent components}, add the component informations
(\cref{thm:p15-w02b-p15-045-additivity}).
\item \emph{For a reduction \(T\)}, compute the induced pair and
\(I'(1:2)\). If \(I(1:2)<\infty\), \(I'(1:2)=I(1:2)\) decides
sufficiency (\cref{thm:p15-w02b-p15-045-reduction}). If
\(I(1:2)=\infty\), equality decides nothing; check instead whether
\(f_1/f_2\) is a function of \(T\) almost everywhere
(\cref{thm:p15-w02b-p15-045-density-criterion}).
\item \emph{Do not reinterpret.} Report \(I(1:2)\) and \(J(1,2)\) as
expected log Bayes factors, not as distances, posterior probabilities
or expected utilities (\cref{sec:p15-w02b-p15-045-distinctions}).
\end{enumerate}
\end{hbmethod}

\section{Relationships and sources}
\label{sec:p15-w02b-p15-045-relations}

\begin{itemize}
\item Finite forms of the directed divergence appear in the
minimal-change updating results of \cref{sec:p15-rat41-dynamics} and in
\cref{def:p16-rat42-kl}. The countable Gibbs inequality with its
safeguards for infinite sums is \cref{thm:p15-rmdl-gibbs}.
\cref{thm:p15-w02b-p15-045-nonnegative} is its measure-theoretic form
under the standing hypotheses.
\item In sequential testing with the score \(\log(f_1/f_2)\) per
observation, the mean score is \(I(1:2)\) under \(H_1\) and \(-I(2:1)\)
under \(H_2\). Wald's lower bounds on expected sample size,
\cref{thm:p15-se-information}, are stated through the two means of such
a score. Relative entropy as the extended mean of a
log score, including the singular case, is
\cref{thm:p15-directional-log-integrability}.
\item The information here is not Fisher's information. Kullback and
Leibler relate the two only in a parametric family, by a second-order
expansion of \(I\) and \(J\) between neighbouring parameter values
\cite[p.~81, eqs.~(2.7)--(2.9)]{KullbackLeibler1951Information}.
\cref{thm:p15-w02b-p15-045-fisher} states that expansion under the
explicit conditions (R1)--(R3) of
\cref{def:p15-w02b-p15-045-regular-family}, which replace the source's
unlisted assumptions. The Fisher-information game of
\cref{sec:p15-huber-information-game} is a different construction.
\item The finite-sample sufficiency result
\cref{thm:p15-r43-sufficiency} is the finite binary case of the
density criterion \cref{thm:p15-w02b-p15-045-density-criterion}.
\item Source: S.~Kullback and R.~A.~Leibler, ``On Information and
Sufficiency'' \cite[pp.~79--86]{KullbackLeibler1951Information}.
Definitions: pp.~79--81; (2.7)--(2.10): pp.~81--82; Lemmas~3.1
and~3.2 and Theorem~3.1: pp.~82--83; Section~4: pp.~83--84; Section~5:
p.~85; Section~6(1), (6.1): p.~85; Section~6(2): pp.~85--86.
\end{itemize}
